% Options for packages loaded elsewhere
\PassOptionsToPackage{unicode}{hyperref}
\PassOptionsToPackage{hyphens}{url}
\documentclass[
  spanish,
]{article}
\usepackage{xcolor}
\usepackage[margin=1in]{geometry}
\usepackage{amsmath,amssymb}
\setcounter{secnumdepth}{-\maxdimen} % remove section numbering
\usepackage{iftex}
\ifPDFTeX
  \usepackage[T1]{fontenc}
  \usepackage[utf8]{inputenc}
  \usepackage{textcomp} % provide euro and other symbols
\else % if luatex or xetex
  \usepackage{unicode-math} % this also loads fontspec
  \defaultfontfeatures{Scale=MatchLowercase}
  \defaultfontfeatures[\rmfamily]{Ligatures=TeX,Scale=1}
\fi
\usepackage{lmodern}
\ifPDFTeX\else
  % xetex/luatex font selection
\fi
% Use upquote if available, for straight quotes in verbatim environments
\IfFileExists{upquote.sty}{\usepackage{upquote}}{}
\IfFileExists{microtype.sty}{% use microtype if available
  \usepackage[]{microtype}
  \UseMicrotypeSet[protrusion]{basicmath} % disable protrusion for tt fonts
}{}
\makeatletter
\@ifundefined{KOMAClassName}{% if non-KOMA class
  \IfFileExists{parskip.sty}{%
    \usepackage{parskip}
  }{% else
    \setlength{\parindent}{0pt}
    \setlength{\parskip}{6pt plus 2pt minus 1pt}}
}{% if KOMA class
  \KOMAoptions{parskip=half}}
\makeatother
\usepackage{longtable,booktabs,array}
\usepackage{calc} % for calculating minipage widths
% Correct order of tables after \paragraph or \subparagraph
\usepackage{etoolbox}
\makeatletter
\patchcmd\longtable{\par}{\if@noskipsec\mbox{}\fi\par}{}{}
\makeatother
% Allow footnotes in longtable head/foot
\IfFileExists{footnotehyper.sty}{\usepackage{footnotehyper}}{\usepackage{footnote}}
\makesavenoteenv{longtable}
\usepackage{graphicx}
\makeatletter
\newsavebox\pandoc@box
\newcommand*\pandocbounded[1]{% scales image to fit in text height/width
  \sbox\pandoc@box{#1}%
  \Gscale@div\@tempa{\textheight}{\dimexpr\ht\pandoc@box+\dp\pandoc@box\relax}%
  \Gscale@div\@tempb{\linewidth}{\wd\pandoc@box}%
  \ifdim\@tempb\p@<\@tempa\p@\let\@tempa\@tempb\fi% select the smaller of both
  \ifdim\@tempa\p@<\p@\scalebox{\@tempa}{\usebox\pandoc@box}%
  \else\usebox{\pandoc@box}%
  \fi%
}
% Set default figure placement to htbp
\def\fps@figure{htbp}
\makeatother
\ifLuaTeX
\usepackage[bidi=basic]{babel}
\else
\usepackage[bidi=default]{babel}
\fi
% get rid of language-specific shorthands (see #6817):
\let\LanguageShortHands\languageshorthands
\def\languageshorthands#1{}
\setlength{\emergencystretch}{3em} % prevent overfull lines
\providecommand{\tightlist}{%
  \setlength{\itemsep}{0pt}\setlength{\parskip}{0pt}}
\usepackage{float}
\usepackage{bookmark}
\IfFileExists{xurl.sty}{\usepackage{xurl}}{} % add URL line breaks if available
\urlstyle{same}
\hypersetup{
  pdftitle={Propuesta de Proyecto: Rendimiento en Powerlifting},
  pdfauthor={Christian Camilo, Luna Garces, Tomas Perez, Tomas Romero, Benjamin Thareau},
  pdflang={es},
  hidelinks,
  pdfcreator={LaTeX via pandoc}}

\title{Propuesta de Proyecto: Rendimiento en Powerlifting}
\usepackage{etoolbox}
\makeatletter
\providecommand{\subtitle}[1]{% add subtitle to \maketitle
  \apptocmd{\@title}{\par {\large #1 \par}}{}{}
}
\makeatother
\subtitle{Metodos Bayesianos -- Primer Semestre 2026}
\author{Christian Camilo, Luna Garces, Tomas Perez, Tomas Romero,
Benjamin Thareau}
\date{}

\begin{document}
\maketitle

\section{Introduccion}\label{introduccion}

El powerlifting es un deporte de fuerza en el que los atletas compiten
principalmente en tres movimientos: sentadilla, press de banca y peso
muerto. A partir de estos levantamientos se obtiene el total levantado,
que corresponde a la suma del mejor intento en cada movimiento. Ademas,
existen puntajes ajustados, como DOTS y Wilks, que permiten comparar el
rendimiento entre atletas con distinto peso corporal.

Este proyecto busca estudiar el rendimiento de atletas de powerlifting
de Estados Unidos utilizando datos reales de OpenPowerlifting. Se busca
analizar como variables como sexo, edad, peso corporal, categoria de
peso y tipo de equipamiento se relacionan con el rendimiento deportivo.
El interes principal es entender que patrones se observan en los datos
y, posteriormente, utilizar modelos bayesianos para responder preguntas
de investigacion.

\section{Datos}\label{datos}

\subsection{Procedencia}\label{procedencia}

Los datos provienen de OpenPowerlifting, una base abierta que recopila
resultados de competencias de powerlifting a nivel internacional. La
base original contiene 3.920.662 registros y 42 variables, incluyendo
informacion del atleta, la competencia, los intentos de levantamiento,
los mejores levantamientos validos, puntajes ajustados y fecha del
evento. Para este proyecto se trabajo con una version filtrada de la
base, restringida a atletas de Estados Unidos (USA).

Es importante notar que cada fila representa un registro competitivo, no
necesariamente una persona unica. Por lo tanto, un mismo atleta puede
aparecer mas de una vez si compitio en distintas fechas o eventos. Esta
estructura es adecuada para analizar el rendimiento observado en
competencias, pero debe considerarse al interpretar resultados, porque
los registros no son completamente independientes a nivel de atleta.

\subsection{Variables involucradas}\label{variables-involucradas}

\begin{longtable}[]{@{}
  >{\raggedright\arraybackslash}p{(\linewidth - 6\tabcolsep) * \real{0.1212}}
  >{\raggedright\arraybackslash}p{(\linewidth - 6\tabcolsep) * \real{0.6061}}
  >{\raggedright\arraybackslash}p{(\linewidth - 6\tabcolsep) * \real{0.1667}}
  >{\raggedright\arraybackslash}p{(\linewidth - 6\tabcolsep) * \real{0.1061}}@{}}
\toprule\noalign{}
\begin{minipage}[b]{\linewidth}\raggedright
Variable
\end{minipage} & \begin{minipage}[b]{\linewidth}\raggedright
Descripcion
\end{minipage} & \begin{minipage}[b]{\linewidth}\raggedright
Tipo
\end{minipage} & \begin{minipage}[b]{\linewidth}\raggedright
Unidad
\end{minipage} \\
\midrule\noalign{}
\endhead
\bottomrule\noalign{}
\endlastfoot
Name & Nombre del atleta. Puede repetirse si el atleta compitio mas de
una vez. & Cualitativa & No aplica \\
Sex & Sexo registrado en la competencia: F o M. & Cualitativa & No
aplica \\
Age & Edad del atleta al momento de la competencia. & Cuantitativa &
Anos \\
AgeClass & Clase de edad usada en la competencia. & Cualitativa & Rango
de edad \\
BodyweightKg & Peso corporal registrado del atleta. & Cuantitativa &
Kg \\
WeightClassKg & Categoria de peso en la que compite el atleta. &
Cualitativa & Kg \\
Equipment & Tipo de equipamiento permitido: Raw, Wraps, Single-ply o
Multi-ply. & Cualitativa & No aplica \\
TipoEquipo & Variable creada para agrupar Raw versus Equipado. &
Cualitativa & No aplica \\
Best3SquatKg & Mejor intento valido de sentadilla dentro de los tres
intentos principales. & Cuantitativa & Kg \\
Best3BenchKg & Mejor intento valido de press de banca dentro de los tres
intentos principales. & Cuantitativa & Kg \\
Best3DeadliftKg & Mejor intento valido de peso muerto dentro de los tres
intentos principales. & Cuantitativa & Kg \\
TotalKg & Suma de los mejores levantamientos validos en sentadilla,
banca y peso muerto. & Cuantitativa & Kg \\
Dots & Puntaje ajustado por peso corporal usado para comparar
rendimiento relativo. & Cuantitativa & Puntos \\
Wilks & Otro puntaje ajustado usado historicamente en powerlifting. &
Cuantitativa & Puntos \\
Date & Fecha de inicio de la competencia. & Fecha & Fecha \\
Year & Ano de la competencia, creado a partir de Date. & Cuantitativa
discreta & Ano \\
Elite & Variable creada para clasificar atletas con Dots \textgreater=
400 como Elite. & Cualitativa & No aplica \\
\end{longtable}

\subsection{Preparacion y limpieza de los
datos}\label{preparacion-y-limpieza-de-los-datos}

Para construir la base final se seleccionaron variables relacionadas con
la identificacion del atleta, caracteristicas fisicas, tipo de
competencia, equipamiento, intentos de levantamiento, mejores marcas,
puntajes de rendimiento y fecha. Luego se eliminaron observaciones con
valores faltantes en las variables centrales del analisis.

Ademas, se aplicaron filtros para trabajar con registros comparables:
atletas de USA, evento SBD, sexo F o M, edades entre 15 y 70 anos, peso
corporal entre 35 y 200 kg, levantamientos positivos, puntaje DOTS mayor
a 10 y competencias desde el ano 2000 en adelante.

\begin{longtable}[]{@{}lr@{}}
\toprule\noalign{}
Indicador & Valor \\
\midrule\noalign{}
\endhead
\bottomrule\noalign{}
\endlastfoot
Registros finales & 10,000 \\
Variables finales & 35 \\
Atletas distintos & 9,481 \\
Primer ano observado & 2000 \\
Ultimo ano observado & 2026 \\
\end{longtable}

\section{Analisis descriptivo}\label{analisis-descriptivo}

\subsection{Medidas de resumen}\label{medidas-de-resumen}

\begin{longtable}[]{@{}lrrrrr@{}}
\toprule\noalign{}
Variable & Media & Mediana & DE & Minimo & Maximo \\
\midrule\noalign{}
\endhead
\bottomrule\noalign{}
\endlastfoot
Edad & 28.77 & 25.00 & 11.48 & 15.00 & 70.00 \\
Peso corporal & 85.76 & 82.20 & 21.94 & 40.64 & 190.50 \\
Mejor sentadilla & 177.85 & 175.00 & 65.64 & 20.00 & 532.97 \\
Mejor banca & 113.75 & 112.50 & 47.72 & 20.41 & 394.63 \\
Mejor peso muerto & 199.19 & 200.00 & 60.37 & 25.00 & 465.00 \\
Total levantado & 490.79 & 490.00 & 167.39 & 75.00 & 1206.56 \\
DOTS & 358.99 & 357.53 & 74.94 & 43.07 & 742.81 \\
Wilks & 356.76 & 355.18 & 74.99 & 43.12 & 734.34 \\
\end{longtable}

La tabla muestra que la base contiene atletas con distintos niveles de
rendimiento, edades y pesos corporales. Esto es importante para el
proyecto, porque permite comparar grupos y estudiar como cambia el
rendimiento segun las caracteristicas fisicas y competitivas.

\subsection{Distribucion de variables
principales}\label{distribucion-de-variables-principales}

\begin{figure}[H]

{\centering \includegraphics[width=0.92\linewidth]{informe_files/figure-latex/histogramas-1} 

}

\caption{Distribucion de edad, peso corporal y puntaje DOTS}\label{fig:histogramas}
\end{figure}

La mayor parte de los registros corresponde a atletas adultos jovenes
(aproximadamente entre 20 y 40 anos). El peso corporal se concentra en
rangos intermedios y el puntaje DOTS muestra una distribucion unimodal,
con mayor densidad en valores medios y una cola derecha asociada a
atletas de rendimiento excepcional. Esto sugiere que la muestra contiene
una mayoria de competidores recreativos o intermedios, junto con un
grupo mas pequeno de atletas de alto nivel.

\subsection{Rendimiento segun sexo}\label{rendimiento-segun-sexo}

\begin{figure}[H]

{\centering \includegraphics[width=0.92\linewidth]{informe_files/figure-latex/boxplot-sexo-1} 

}

\caption{Distribucion del puntaje DOTS segun sexo}\label{fig:boxplot-sexo}
\end{figure}

En la muestra, los hombres presentan una mediana de DOTS algo mayor que
las mujeres, el rendimiento tambien puede depender de la edad, el peso
corporal, la categoria de peso y el tipo de equipamiento.

\begin{figure}[H]

{\centering \includegraphics[width=0.92\linewidth]{informe_files/figure-latex/densidad-sexo-1} 

}

\caption{Densidad del puntaje DOTS segun sexo}\label{fig:densidad-sexo}
\end{figure}

El grafico de densidad confirma que las distribuciones no son identicas.
La curva de hombres aparece desplazada hacia valores mas altos, aunque
existe un solapamiento importante entre ambos grupos.

\subsection{Rendimiento segun tipo de
equipamiento}\label{rendimiento-segun-tipo-de-equipamiento}

\begin{figure}[H]

{\centering \includegraphics[width=0.92\linewidth]{informe_files/figure-latex/boxplot-equip-1} 

}

\caption{Puntaje DOTS segun tipo de equipamiento}\label{fig:boxplot-equip}
\end{figure}

El puntaje DOTS cambia segun el tipo de equipamiento. Las categorias
equipadas, especialmente \emph{Single-ply} y \emph{Multi-ply}, tienden a
presentar mayor dispersion y valores altos.

\begin{figure}[H]

{\centering \includegraphics[width=0.92\linewidth]{informe_files/figure-latex/boxplot-deadlift-1} 

}

\caption{Peso muerto en clase 90 kg segun equipamiento}\label{fig:boxplot-deadlift}
\end{figure}

Al fijar la clase de peso en 90 kg se reduce parte de la variabilidad
asociada al peso corporal. Las diferencias entre equipamientos permiten
observar de forma mas directa si el uso de soporte se asocia con mejores
marcas de peso muerto.

\subsection{Relacion entre Wilks y total
levantado}\label{relacion-entre-wilks-y-total-levantado}

\begin{figure}[H]

{\centering \includegraphics[width=0.92\linewidth]{informe_files/figure-latex/wilks-total-1} 

}

\caption{Relacion entre puntaje Wilks y total levantado segun sexo}\label{fig:wilks-total}
\end{figure}

Existe una relacion positiva clara entre Wilks y el total levantado.
Esto es esperable, porque Wilks resume parte del rendimiento ajustando
por peso corporal. Ademas, se observa separacion entre hombres y mujeres
en el total absoluto levantado.

\subsection{Promedio de levantamientos segun
edad}\label{promedio-de-levantamientos-segun-edad}

\begin{figure}[H]

{\centering \includegraphics[width=0.92\linewidth]{informe_files/figure-latex/levantamientos-edad-1} 

}

\caption{Promedio de sentadilla, banca y peso muerto segun clase de edad}\label{fig:levantamientos-edad}
\end{figure}

El promedio levantado aumenta desde las edades mas jovenes hasta los 30
y luego tiende a disminuir en edades mayores. El peso muerto y la
sentadilla suelen presentar promedios mas altos que la banca. Esto
sugiere que la edad no tiene un efecto lineal simple sobre el
rendimiento.

\subsection{Relacion entre edad y
DOTS}\label{relacion-entre-edad-y-dots}

\begin{figure}[H]

{\centering \includegraphics[width=0.92\linewidth]{informe_files/figure-latex/dispersion-dots-1} 

}

\caption{Relacion entre edad y puntaje DOTS segun sexo}\label{fig:dispersion-dots}
\end{figure}

La relacion entre edad y DOTS no parece completamente lineal. El
rendimiento aumenta en las primeras etapas, alcanza niveles altos en
adultos jovenes y luego muestra una tendencia a disminuir en edades
mayores.

\subsection{Relacion entre peso corporal y
DOTS}\label{relacion-entre-peso-corporal-y-dots}

\begin{figure}[H]

{\centering \includegraphics[width=0.92\linewidth]{informe_files/figure-latex/dispersion-peso-1} 

}

\caption{Relacion entre peso corporal y puntaje DOTS segun sexo}\label{fig:dispersion-peso}
\end{figure}

Como DOTS ya ajusta por peso corporal, no se espera una relacion lineal
fuerte como ocurriria con el total absoluto levantado. El suavizamiento
sugiere que el rendimiento relativo puede variar entre rangos de peso,
con zonas donde el puntaje promedio se estabiliza y otras donde
disminuye. Esto motiva modelar la relacion de forma flexible en vez de
imponer una tendencia lineal simple.

REVISAR RESPUESTA

\subsection{DOTS por sexo y clase de
edad}\label{dots-por-sexo-y-clase-de-edad}

\begin{figure}[H]

{\centering \includegraphics[width=0.92\linewidth]{informe_files/figure-latex/dots-edad-sexo-1} 

}

\caption{Puntaje DOTS por sexo y clase de edad}\label{fig:dots-edad-sexo}
\end{figure}

El grafico por clase de edad muestra que las diferencias entre sexos se
mantienen en varias categorias etarias, aunque la magnitud de la
diferencia cambia segun el grupo. Tambien se aprecia que algunas clases
de edad presentan mayor variabilidad, posiblemente por diferencias en
tamano muestral o por la mezcla de atletas con trayectorias deportivas
distintas.

\subsection{DOTS promedio por rangos de edad y
peso}\label{dots-promedio-por-rangos-de-edad-y-peso}

\begin{figure}[H]

{\centering \includegraphics[width=0.92\linewidth]{informe_files/figure-latex/heatmap-dots-1} 

}

\caption{DOTS promedio por rangos de edad y peso corporal}\label{fig:heatmap-dots}
\end{figure}

El mapa de calor resume el rendimiento relativo promedio (medido en
DOTS) según rangos de edad y peso corporal. Aunque la métrica DOTS está
diseñada en teoría para equilibrar la ventaja del peso, el gráfico
revela que los puntajes más altos no se distribuyen uniformemente,
concentrándose fuertemente en el rango de 20 a 40 años y las categorías
más pesadas (sobre los 120 kg). El gráfico muestra la interacción entre
ambas variables: el declive del rendimiento asociado al envejecimiento
no ocurre con la misma intensidad ni al mismo ritmo en todas las
categorías de peso. Las celdas con menos registros fueron excluidas para
garantizar que estas tendencias representen comportamientos
poblacionales estables.

\subsection{Correlacion entre variables
numericas}\label{correlacion-entre-variables-numericas}

\begin{figure}[H]

{\centering \includegraphics[width=0.92\linewidth]{informe_files/figure-latex/matriz-correlacion-1} 

}

\caption{Matriz de correlacion entre variables numericas principales}\label{fig:matriz-correlacion}
\end{figure}

La matriz muestra correlaciones altas entre los tres levantamientos y el
total, lo que es esperable porque el total se construye directamente a
partir de ellos. DOTS y Wilks tambien presentan una correlacion muy
alta, ya que ambos son puntajes ajustados por peso corporal y buscan
medir rendimiento relativo. En cambio, la edad suele tener asociaciones
mas debiles con las variables de fuerza, lo que sugiere que su efecto
puede depender de otros factores.

\subsection{Evolucion de competidores en el
tiempo}\label{evolucion-de-competidores-en-el-tiempo}

\begin{figure}[H]

{\centering \includegraphics[width=0.92\linewidth]{informe_files/figure-latex/evolucion-temporal-1} 

}

\caption{Evolucion de participantes por año segun tipo de equipamiento}\label{fig:evolucion-temporal}
\end{figure}

La evolucion temporal muestra cambios importantes en la participacion y
en la composicion por tipo de equipamiento. En particular, los registros
\emph{Raw} tienden a ganar protagonismo en los años recientes, mientras
que las categorias equipadas quedan relativamente más concentradas en
periodos o grupos especificos.

\section{Preguntas de investigacion}\label{preguntas-de-investigacion}

Las siguientes preguntas surgen del interes del equipo por comprender
distintas dimensiones del rendimiento en powerlifting competitivo,
considerando tanto factores fisicos como el rol del equipamiento y la
evolucion del deporte en el tiempo.

\textbf{Pregunta 1:} Existe una diferencia significativa en la media de
los kilogramos levantados en peso muerto entre atletas \emph{Raw} y
aquellos con equipamiento dentro de una misma clase de peso?

\begin{quote}
\emph{Motivacion:} Esta pregunta permite aislar parcialmente el efecto
del equipamiento al comparar atletas dentro de una misma categoria de
peso. El interes bayesiano seria estimar la diferencia de medias entre
grupos y cuantificar la incertidumbre asociada a esa diferencia.
\end{quote}

\textbf{Pregunta 2:} Como ha evolucionado la cantidad total de
competidores a lo largo del tiempo dependiendo de si utilizan o no
equipamiento de soporte?

\begin{quote}
\emph{Motivacion:} La popularidad relativa de las categorias \emph{Raw}
y equipadas puede haber cambiado en el tiempo. Modelar esta evolucion
permite estudiar tendencias temporales y comparar tasas de crecimiento
entre tipos de competencia.
\end{quote}

\textbf{Pregunta 3:} Hay diferencias significativas en el rendimiento
(puntaje DOTS) entre hombres y mujeres en USA, controlando por edad y
peso corporal?

\begin{quote}
\emph{Motivacion:} La comparacion directa por sexo puede estar
confundida por diferencias de edad, peso corporal y composicion de la
muestra. Un modelo con estas variables fijas permitiria evaluar
diferencias en rendimiento.
\end{quote}

\textbf{Pregunta 4:} Cuales caracteristicas del atleta actuan como
factores predictivos para alcanzar el nivel de elite en USA?

\begin{quote}
\emph{Motivacion:} Clasificar atletas como elite o no elite permite
formular un problema predictivo. Variables como edad, sexo, peso
corporal, equipamiento y levantamientos pueden aportar informacion sobre
la probabilidad de alcanzar un nivel alto de rendimiento.
\end{quote}

\textbf{Pregunta 5:} Cual es la relacion entre la edad, el peso corporal
y el puntaje DOTS en atletas de powerlifting en USA?

\begin{quote}
\emph{Motivacion:} Los gráficos descriptivos muestran que la relación
entre edad y rendimiento no es lineal: el puntaje DOTS aumenta en
atletas jóvenes y tiende a disminuir en edades mayores. Se utilizara un
enfoque bayesiano, el cual permite incorporar incertidumbre en la
estimación de los coeficientes y es especialmente útil en zonas con
menos datos, como las edades extremas.
\end{quote}

\section{Discusión preliminar}\label{discusiuxf3n-preliminar}

A partir del analisis descriptivo realizado, es posible identificar
algunas tendencias generales que orientan las preguntas de investigacion
planteadas. Primero, el rendimiento medido mediante DOTS presenta
diferencias entre sexos. Para una comparación justa, se deben incorporar
diferentes parámetros de efecto para cada sexo, controlando además por
edad y peso corporal, en lugar de basarse únicamente en diferencias
visuales de mediana.

Segundo, el equipamiento parece estar relacionado con cambios en el
rendimiento y con patrones distintos de participacion. Las categorias
equipadas muestran mayor dispersion en algunos graficos, mientras que
los registros \emph{Raw} parecen tener una presencia creciente en el
tiempo. Estas observaciones justifican preguntas especificas sobre el
efecto del equipamiento y sobre la evolucion temporal de las categorias.

Tercero, la edad y el peso corporal no se relacionan con el rendimiento
de forma puramente lineal. El rendimiento promedio tiende a aumentar
desde edades tempranas, estabilizarse en adultos jovenes y disminuir en
edades mayores. Ademas, el mapa de calor sugiere que ciertas
combinaciones de edad y peso concentran mejores puntajes promedio. Por
lo tanto, los modelos posteriores deberian considerar posibles efectos
no lineales e interacciones.

Finalmente, el analisis confirma que DOTS y Wilks capturan informacion
similar sobre rendimiento relativo, mientras que el total levantado y
los tres levantamientos principales estan fuertemente relacionados entre
si. Esto ayuda a decidir que variables usar como respuesta y cuales
utilizar como predictores, evitando redundancias innecesarias en los
modelos bayesianos.

\section{Referencias}\label{referencias}

\begin{itemize}
\tightlist
\item
  OpenPowerlifting. (s.f.). \emph{Open Powerlifting Database}.
  Recuperado de \url{https://www.openpowerlifting.org}
\end{itemize}

\end{document}
